\section{Acoplamiento reversible y transporte no estacionario}
\label{nuclear:10}
\subsection{Extensión reversible del balance entre dos formas}

Sean \(B,C:V\to V'\) isomorfismos que conservan una forma no degenerada \(b\) entre dos realizaciones:



\[
B^*b'B=C^*b'C=b.
\]



La forma puede ser una métrica positiva o una forma alternante. En espacios reales, \(*\) denota transposición respecto de las coordenadas elegidas; la igualdad expresa el transporte de la forma entre las realizaciones especificadas. En espacios de Hilbert, los operadores y sus inversos son acotados.

Con \(s=\sqrt8\) y \(Q_V=\frac13\left(\begin{smallmatrix}sI&-I\\I&sI\end{smallmatrix}\right)\), definimos



\begin{equation}
\boxed{
\mathcal U(B,C)=Q_{V'}^*\operatorname{diag}(B,C)Q_V
=\frac19\begin{pmatrix}8B+C&\sqrt8(C-B)\\\sqrt8(C-B)&B+8C\end{pmatrix}.}
\label{nuc10:eq:1}
\end{equation}



Escribimos \(T=(8B+C)/9\), \(D=\sqrt8(C-B)/9\) y \(R=(B+8C)/9\). La actualización completa es



\begin{equation}
\binom{x'}{m'}=\begin{pmatrix}T&D\\D&R\end{pmatrix}\binom{x}{m}.
\label{nuc10:eq:2}
\end{equation}



\begin{theorem}
La aplicación \eqref{nuc10:eq:1} es reversible y conserva la forma \(b\oplus b\):



\begin{equation}
\mathcal U(B,C)^*(b'\oplus b')\mathcal U(B,C)=b\oplus b,
\quad\mathcal U(B,C)^{-1}=\mathcal U(B^{-1},C^{-1}).
\label{nuc10:eq:3}
\end{equation}
\end{theorem}



\begin{proof}
Se cumple \(Q^*Q=I\) por \(8+1=9\). Sus bloques escalares hacen que \(Q\) preserve \(b\oplus b\) para cualquier forma \(b\). La aplicación diagonal conserva las formas por hipótesis. La composición prueba \eqref{nuc10:eq:3}, y la inversión de los tres factores da la inversa declarada.
\end{proof}

Para memoria entrante \(m=0\), se recupera el balance \(b=T^*b'T+D^*b'D\). La segunda columna determina la respuesta a memoria entrante arbitraria. La conservación de normas positivas y la conservación de área orientada son las especializaciones respectivas de \eqref{nuc10:eq:3}.

\subsubsection{Realización colectiva en las nueve copias}

En \(V^9\), sea \(W_E=\operatorname{diag}(I,\ldots,I,E)\), con \(E\) en la novena posición. Sean \(A(a,b)=(a/\sqrt8,\ldots,a/\sqrt8,b)\) y \(L=AQ\). Entonces \(L(x,0)=Jx\) y



\begin{equation}
W_E L=L\mathcal U(I,E).
\label{nuc10:eq:4}
\end{equation}



\begin{proof}
\(W_EA=A\operatorname{diag}(I,E)\); multiplicar por \(Q\) y usar \(AQ=L\) da \eqref{nuc10:eq:4}. La ortogonalidad de la media de las ocho primeras posiciones con sus siete diferencias muestra además que el complemento queda fijo bajo \(W_E\).
\end{proof}

El paso cíclico original es \(S_E=P_{\rm cyc}W_E\). Por ello \(S_EL=(P_{\rm cyc}L)U(I,E)\), con coordenatizaciones de entrada y salida distintas. La reducción colectiva conserva el patrón correspondiente: \(S_E^9=I_9\otimes E\), mientras \(W_E^9=\operatorname{diag}(I_8,E^9)\). Cuando un lector utiliza las nueve etiquetas individuales, éstas se conservan en su representación original.

\subsection{Composición exacta y reentrada de memoria}

Para \(B_1,C_1:V_0\to V_1\) y \(B_2,C_2:V_1\to V_2\),



\begin{equation}
\boxed{\mathcal U(B_2,C_2)\mathcal U(B_1,C_1)
=\mathcal U(B_2B_1,C_2C_1).}
\label{nuc10:eq:5}
\end{equation}



\begin{proof}
En \eqref{nuc10:eq:1} se cancelan \(Q_{V_1}Q_{V_1}^*\), y los dos bloques diagonales se multiplican en el orden indicado. La inducción extiende la fórmula a cualquier número finito de transportes, conservando su orden.
\end{proof}

El bloque de lectura contiene la identidad



\begin{equation}
\boxed{T_2T_1+D_2D_1=\frac{8B_2B_1+C_2C_1}{9}.}
\label{nuc10:eq:6}
\end{equation}



Desde \((x,0)\), la primera etapa produce \((T_1x,D_1x)\). La segunda lectura es \(T_2T_1x+D_2D_1x\). El segundo término es exactamente la contribución de la memoria reintroducida. Al almacenar \(D_1x\) fuera del siguiente acoplamiento y continuar sólo con \(T_1x\) se obtiene el procedimiento de registros separados, cuya lectura es \(T_2T_1x\). La igualdad \eqref{nuc10:eq:6} precisa la diferencia entre ambos procedimientos.

La realización \(U\) conserva el estado del acoplamiento y los productos ordenados \(B_N\cdots B_1\), \(C_N\cdots C_1\). La conservación de todos los prefijos incluye además el registro de emisiones de HMT, que mantiene cada factor junto con su orden.

\subsubsection{Ecuación covariante de incremento}

Para el registro separado \(x_{n+1}=T_nx_n\), \(m_n=D_nx_n\), se tiene



\begin{equation}
x_{n+1}-B_nx_n=m_n/\sqrt8.
\label{nuc10:eq:7}
\end{equation}



Con \(G_0=I\) y \(G_{n+1}=B_nG_n\), la telescopía da



\begin{equation}
x_0=G_N^{-1}x_N-\frac1{\sqrt8}\sum_{n<N}G_{n+1}^{-1}m_n.
\label{nuc10:eq:8}
\end{equation}



La memoria es el incremento respecto del transporte entre formas \(B_n\), con la normalización declarada. La igualdad \eqref{nuc10:eq:8} es una reconstrucción finita algebraica; su paso a una serie infinita requiere convergencia. El teorema siguiente construye la inversa estable mediante el límite de operadores positivos.

\subsection{Memoria de una sucesión no estacionaria}

En una realización de Hilbert \(H\), sean \(C_n\) unitarios, con su orden de composición declarado. Escribimos



\[
T_n=(8I+C_n)/9,\quad D_n=\sqrt8(C_n-I)/9,
\quad P_0=I,\quad P_{n+1}=T_nP_n,\quad A_N=P_N^*P_N.
\]



El índice \(N\) cuenta compresiones con almacenamiento separado, a diferencia del transporte completo con reentrada. Para métricas variables con \(B_n\) isomorfismos isométricos, la conjugación por \(G_n\) lleva los transportes a un Hilbert fijo y produce este mismo caso, con \(\widetilde C_n=G_{n+1}^{-1}C_nG_n\). La afirmación utiliza las métricas de las fibras y las isometrías declaradas entre ellas.

\begin{theorem}[Terminal positivo y recuperación no estacionaria]
\label{nuclear:terminal-no-estacionario}
Existe un operador positivo \(0\leq A_\infty\leq I\), límite fuerte de \(A_N\), y



\begin{equation}
\boxed{I=A_\infty+\sum_{n\ge0}P_n^*D_n^*D_nP_n.}
\label{nuc10:eq:9}
\end{equation}



La aplicación



\begin{equation}
\boxed{\mathcal Vv=(A_\infty^{1/2}v,(D_nP_nv)_{n\ge0})}
\label{nuc10:eq:10}
\end{equation}



es isométrica y tiene inversa izquierda estable



\begin{equation}
\mathcal V^*(a,(m_n))=A_\infty^{1/2}a+
\sum_{n\ge0}P_n^*D_n^*m_n,\qquad\|\mathcal V^*\|=1.
\label{nuc10:eq:11}
\end{equation}



\end{theorem}
\begin{proof}
La identidad local \(T_n^*T_n+D_n^*D_n=I\) da
\(A_N-A_{N+1}=P_N^*D_N^*D_NP_N\geq0\). Las formas cuadráticas de
la sucesión positiva decreciente convergen a la de un operador positivo
\(A_\infty\). Para \(B_N=A_N-A_\infty\), la desigualdad
\(0\leq B_N^2\leq B_N\) convierte esa convergencia en fuerte.
La telescopía finita y el límite prueban \eqref{nuc10:eq:9}. Evaluar sobre \(v\)
prueba \eqref{nuc10:eq:10}. Su adjunto es \eqref{nuc10:eq:11}; la serie converge en norma porque
las truncaciones de un registro cuadrado-sumable convergen y el adjunto es acotado.
\end{proof}

La reconstrucción con el terminal límite y \(N\) registros tiene error
exacto \((A_N-A_\infty)v\). El vector \(A_\infty^{1/2}v\) es la
realización positiva canónica del terminal de norma. Su existencia procede
del límite de Gram, con independencia de la convergencia vectorial de \(P_Nv\).

\subsection{Tres propiedades diferentes: agotamiento, distinción y estabilidad}

Sea \(Mv=(D_nP_nv)_n\). De \eqref{nuc10:eq:9}, \(M^*M=I-A_\infty\). Por tanto:

\begin{enumerate}
\item Toda la norma de \(v\) pasa al registro exactamente cuando \(A_\infty v=0\).
\item El registro distingue cada estado exactamente cuando \(\ker(I-A_\infty)=\{0\}\).
\item La recuperación desde memoria sola tiene inversa acotada exactamente cuando
\(I-A_\infty\geq\varepsilon I\) para algún \(\varepsilon>0\).
En ese caso una inversa izquierda es \((I-A_\infty)^{-1}M^*\).
\end{enumerate}

Además,



\begin{equation}
\boxed{\ker M=\bigcap_{n\ge0}\operatorname{Fix}(C_n).}
\label{nuc10:eq:12}
\end{equation}



\begin{proof}
Si \(Mv=0\), entonces \(D_0v=0\), de donde \(C_0v=v\)
y \(P_1v=v\). Inductivamente, \(D_nP_nv=0\) da \(C_nv=v\) y
\(P_{n+1}v=v\). El recíproco es inmediato. Las tres equivalencias se
deducen del Gram \(M^*M\) y de la caracterización de los operadores
acotados inferiormente.
\end{proof}

\subsubsection{Ejemplo exacto que separa las tres propiedades}

En el ejemplo algebraico \(C_0=-I\) y \(C_n=I\) para \(n\geq1\),



\begin{equation}
A_\infty=\frac{49}{81}I,\qquad M^*M=\frac{32}{81}I,
\qquad m_0=-\frac{2\sqrt8}{9}v,
\qquad v=-\frac9{2\sqrt8}m_0.
\label{nuc10:eq:13}
\end{equation}



La memoria contiene \(32/81\) de la norma al cuadrado y permite recuperar el vector completo. El terminal retiene \(49/81\) y es un operador positivo distinto de un proyector. La recuperabilidad depende del núcleo y del condicionamiento del lector. Lectura y memoria constituyen componentes correlacionadas de una isometría.

Este ejemplo separa las tres consecuencias lógicas del teorema. La
aplicación a una trayectoria física conserva además la selección y la
procedencia efectiva de sus \(C_n\) mediante el TPK.

\subsection{Criterio cuantitativo para transferencia completa}

Para \(P_{j,n}=T_{j-1}\cdots T_n\) y \(P_{n,n}=I\), una ventana de longitud \(L\) tiene Gram



\[
G_{n,L}=\sum_{j=n}^{n+L-1}P_{j,n}^*D_j^*D_jP_{j,n}
=I-P_{n+L,n}^*P_{n+L,n}.
\]



Si \(G_{n,L}\geq\kappa I\), con \(\kappa>0\), en cada bloque consecutivo de longitud \(L\), entonces



\begin{equation}
\|P_N\|^2\le(1-\kappa)^{\lfloor N/L\rfloor},\qquad A_\infty=0.
\label{nuc10:eq:14}
\end{equation}



Cada bloque contrae la norma al cuadrado por \(1-\kappa\); la composición y la contractividad de los pasos restantes prueban \eqref{nuc10:eq:14}.

Un criterio suficiente anterior a los productos es



\begin{equation}
\sum_{j=n}^{n+L-1}D_j^*D_j\ge\alpha I
\quad\Longrightarrow\quad
G_{n,L}\ge\frac{\alpha}{[1+2(L-1)/9]^2}I.
\label{nuc10:eq:15}
\end{equation}



\begin{proof}
Para \(y_j=P_{j,n}v\), se tiene
\(y_{j+1}-y_j=D_jy_j/\sqrt8\) y \(\|D_j\|\leq2\sqrt8/9\). Por ello
\[
\|D_jv\|\leq\|D_jy_j\|+\frac29\sum_{i=n}^{j-1}\|D_iy_i\|.
\]
La matriz triangular de sumas parciales tiene norma a lo sumo \(L-1\).
Aplicar la norma \(\ell^2\) de la ventana y la hipótesis de Gram prueba \eqref{nuc10:eq:15}.

\end{proof}

La ausencia de un vector fijo común da distinción, mientras \eqref{nuc10:eq:15} añade una separación uniforme cuantitativa. Ninguna de las dos condiciones se presume por mencionar simetrías: se aplica a los transportes efectivos de la secuencia.

\subsection{Realización incidencial íntegra y cambios de régimen}

La sección~\ref{nuc05:isometria} construye, para el registro de doce
coordenadas, \(Fk=(\mu(k),\mathcal IP_0k/6)\), con inversa y
producto interno ponderado en \(Y=\RR\oplus\mathcal I(\mathbf1^\perp)\).
La transformación \(F\oplus F\) lleva el acoplamiento completo a



\begin{equation}
\mathcal U_Y=(F'\oplus F')\mathcal U(B,C)(F^{-1}\oplus F^{-1})
=\mathcal U(F'BF^{-1},F'CF^{-1}).
\label{nuc10:eq:16}
\end{equation}



Esta identidad conserva la segunda columna, la reentrada y la composición.
Para formas variables se transporta también la forma mediante \(F\).
En la realización real de las doce fibras \(A_2\), se usa
\(F\otimes\operatorname{id}_{V_{A_2}}\), con
\(V_{A_2}=A_2\otimes_{\ZZ}\RR\). Su dominio es el espacio real de
fibras; la preservación del retículo entero corresponde al transporte
reticular ya demostrado. La aplicación \(F\) recupera \(K\) y
conserva las restantes etiquetas genealógicas en el estado completo.

La versión positiva de \eqref{nuc10:eq:9}–\eqref{nuc10:eq:15} utiliza transportes unitarios en las
métricas declaradas. La versión de forma alternante de \eqref{nuc10:eq:1}–\eqref{nuc10:eq:8} admite
los tres regímenes de Euler. Con
\(E=F_{\rm av}^2=\left(\begin{smallmatrix}1&1\\1&2\end{smallmatrix}\right)\),
se obtiene \(\det T=89/81\), \(\det D=-8/81\). Las sumas
orientadas finitas telescopan y los terminales de área crecen como
\((89/81)^N\). El límite positivo de Hilbert utiliza las hipótesis
unitarias específicas del teorema~\ref{nuclear:terminal-no-estacionario}.

El balance reúne formas y dominios, evolución reversible con reentrada
y recuperación ilimitada en las realizaciones positivas. Las
aplicaciones físicas siguientes conservan los lectores, conexiones y
escalas de esa composición operatoria.
